\section{Related Work}
\label{sec:related-work}

\subsection{Memory Representations and Management in LLM Agents}

Memory is commonly treated as a core component of LLM agents alongside planning and tool use \cite{zhao2023llm_agent_survey,wang2024survey_autonomous_agents}. Existing systems differ in both their representational substrate and their management policy \cite{zhang2024survey_memory_llm_agents}. Explicit memory stores text, structured records, summaries, or retrieved documents. These forms are inspectable and can preserve source-level evidence, but they occupy context space when exposed to the model. Latent memory instead retains internal activations, learned states, or memory tokens. It is more tightly coupled to model computation, but its content is harder to inspect and verify.

This distinction does not by itself define a complete memory system. A long-horizon agent also needs rules for writing, selecting, updating, replacing, and clearing stored information. Our work focuses on these lifecycle operations for latent memories generated by an existing model. The contribution is therefore a management layer over generated latent state, rather than a new taxonomy of agent memory.

\subsection{Recurrent and Retrieved Latent State}

Several architectures preserve hidden state beyond a local segment. Transformer-XL reuses cached hidden states through segment-level recurrence \cite{dai2019transformerxl}, while Compressive Transformer retains older states in a compressed form \cite{rae2019compressive}. Recurrent Memory Transformer passes dedicated memory tokens between segments \cite{bulatov2022rmt}, and Infini-attention integrates compressive memory into the attention mechanism \cite{munkhdalai2024infini}. These methods embed persistence into the model's recurrent computation and typically require architecture-specific training.

Other approaches retrieve stored internal representations. Memorizing Transformers performs approximate nearest-neighbour lookup over a non-differentiable store of internal key-value pairs \cite{wu2022memorizing}. CAMELoT attaches a training-free consolidated associative memory to a frozen language model \cite{he2024camelot}. MEMORYLLM maintains a self-updatable latent memory pool \cite{wang2024memoryllm}, and M+ combines latent memory with a trained retriever for longer retention \cite{wang2025mplus}. Titans instead learns a neural long-term memory that is updated at test time \cite{behrouz2025titans}.

Our method shares the goal of persistent internal state but differs in its integration point. It does not add a new recurrent architecture or train a new memory module. It manages latent sequences already produced by the frozen MemGen pathway within a bounded session-owned store.

\subsection{Generated Latent Memory in MemGen}

MemGen dynamically generates machine-native memory during reasoning \cite{zhang2025memgen}. A memory Trigger decides when augmentation is needed, and a Weaver constructs a latent token sequence that supports the Reasoner. This design makes memory generation conditional on the current reasoning state.

We build on this mechanism rather than replace it. Our question begins after a Weaver output has been generated. We ask who owns that memory, how long it should remain available, and which memories a later query should retrieve. We also define how a bounded store updates or replaces its contents. The proposed bank makes these lifecycle semantics explicit. It resets state between sessions and blocks writes during frozen-bank question answering. This operational focus distinguishes our work from methods whose main contribution is the generation or training of latent memory itself.

\subsection{Explicit Memory and Long-Horizon Evaluation}

Explicit-memory systems provide an important comparison because they retain readable evidence. Retrieval-augmented generation supplies selected documents to a parametric model \cite{lewis2020rag}, while MemLLM trains a model to interact with an explicit read-write memory \cite{modarressi2024memllm}. MemGPT manages multiple memory tiers as virtual context for long-running interactions \cite{packer2023memgpt}. In our EventQA evaluation, same-model summaries, BM25 retrieval, and a matched-token-budget condition test whether readable memory or additional query-time text explains the observed effect.

Long-horizon memory benchmarks also impose different evidence contracts. LongMemEval tests sustained conversational memory, including temporal reasoning, knowledge updates, and abstention \cite{wu2024longmemeval}. LoCoMo evaluates question answering, summarization, and dialogue generation over long conversations \cite{maharana2024locomo}. MemoryAgentBench instead evaluates several memory competencies through incremental interactions and includes EventQA as an accurate-retrieval task \cite{hu2026evaluating}.
